\begin{table}[!htbp]
\centering
\begingroup\singlespacing\fontsize{8.6}{10.58}\selectfont
\setlength{\tabcolsep}{3pt}
\caption{重复使用留出集的探索性配对比较}\label{tab:S10}
\begin{tabular}{@{}>{\raggedright\arraybackslash}p{0.24\linewidth}>{\raggedright\arraybackslash}p{0.075\linewidth}>{\raggedright\arraybackslash}p{0.13\linewidth}>{\raggedright\arraybackslash}p{0.22\linewidth}>{\raggedright\arraybackslash}p{0.1\linewidth}>{\raggedright\arraybackslash}p{0.1\linewidth}@{}}
\toprule
比较 & 结局 & $\Delta$RMSE & 95\% CI & $p$ (swap) & $p$ (Holm) \\
\midrule
候选 $-$ 均值 & NA & $-0.003408$ & [$-0.0067$, $-0.0004$] & 0.034 & 0.168 \\
候选 $-$ 均值 & PA & $-0.005163$ & [$-0.0082$, $-0.0021$] & 0.002 & 0.011 \\
候选 $-$ 均值 & LS & $-0.002525$ & [$-0.0064$, 0.0014] & 0.204 & 0.817 \\
候选 $-$ 原 RF & NA & $-0.000023$ & [$-0.0022$, 0.0022] & 0.986 & 1.000 \\
候选 $-$ 原 RF & PA & $-0.000396$ & [$-0.0025$, 0.0016] & 0.712 & 1.000 \\
候选 $-$ 原 RF & LS & $-0.000091$ & [$-0.0029$, 0.0027] & 0.948 & 1.000 \\
\bottomrule
\end{tabular}
\tabnote{NA 为消极情绪，PA 为积极情绪，LS 为生活满意度。$n=104$。差值为候选模型减去比较模型。逐点区间来自 5,000 次配对 bootstrap，检验采用 10,000 次 MSE 预测交换，Holm 覆盖六行。后三行比较探索性候选模型与主分析共享 RF。这些比较属于重复使用留出集的探索性分析。}
\endgroup
\end{table}
